\section{What Crust is}

Crust is a Rust subset that compiles through a C compiler. A single
translation unit may hold C, Rust and rpython at once, and after lowering
there is no boundary between them: everything becomes C, and the C compiler
sees one program.

\begin{verbatim}
#include "schemes.py"        /* rpython: string and list work         */
int printf(const char *, ...);
struct Frames { words: [u64; 8], free: i64 }   /* Rust: fixed layouts  */
impl Frames {
    fn alloc(&mut self) -> i64 { ... }
}
int main(void) { ... }                          /* C: entry point      */
\end{verbatim}

The point of mixing three languages is that each is short at something
different. rpython has lists, strings and a `\%'-formatting runtime, so URL
parsing and table building are a few lines there and a rewrite in the Rust
subset. Rust has fixed layouts, traits and generics with no allocator
underneath, so a bitmap allocator or a scheduler is natural. C is the entry
point and the linkage surface.

Nothing is copied at the boundaries. \texttt{tools/py2c.py} lowers a typed
rpython list to three words --- \texttt{\{ int* data; long len; long cap; \}}
--- and \texttt{PyList<T>} in Crust's bundled core has exactly that layout, so
a list rpython built is walked directly from Rust with a pointer cast and no
conversion.

\subsection{What is in the subset}

Structs, enums (including data-carrying ones, lowered to tagged unions),
traits with default methods and associated consts and types, generics by
monomorphisation, closures lifted to static functions, slices, \texttt{match}
with bindings, \texttt{?}, \texttt{impl} blocks, operator sugar for the
bundled containers, and a small core: \texttt{Vec<T>}, \texttt{Box<T>},
\texttt{String}, \texttt{Formatter} with \texttt{write!} and \texttt{format!},
\texttt{Rc}/\texttt{Arc}, \texttt{VecDeque<T>}, and the \texttt{core} free
functions real code reaches for (\texttt{slice::from\_raw\_parts},
\texttt{ptr::null\_mut}, \texttt{cmp::min}).

Generics are monomorphised, so a trait call is a direct call and an
instantiation costs what a hand-written version would --- the
\texttt{traits} and \texttt{generics} benchmarks below are there to check that
claim rather than assert it.

\subsection{What is deliberately absent}

There is no borrow checker, no \texttt{dyn Trait}, no iterator protocol, and
no \texttt{Drop}. Each of those is a decision rather than an omission waiting
to be filled in:

\begin{itemize}
  \item \textbf{No \texttt{Drop}} means every allocating type has an explicit
    \texttt{free\_buf}. Scope exit cannot run code, so pretending otherwise
    would leak silently.
  \item \textbf{No iterator protocol}, because rpython already has one and its
    output is a plain struct. Iteration-heavy code is written on that side.
  \item \textbf{\texttt{Mutex<T>} does not lock.} It exists so source written
    against \texttt{std} parses. Crust has no threads --- no spawn, no atomics,
    no memory model --- so there is nothing for a lock to protect against, and
    a lock that does nothing is consistent with the rest of the model rather
    than a hole in it. The moment real concurrency exists it becomes dangerous
    and must be replaced, not extended.
\end{itemize}

\section{CrustOS, and what it borrows from Redox}

CrustOS is a small operating system prototype assembled from the parts of
\href{https://redox-os.org}{Redox} that Crust can compile, plus the parts
supplied here. \texttt{tools/crustos.py} clones the Redox kernel and relibc,
compiles the compatible subset, and links it with a scheme layer written in
rpython and a scheduler written in Rust.

\begin{verbatim}
python3 tools/crustos.py fetch    # clone redox-os/kernel and relibc
python3 tools/crustos.py run      # compile the subset, link, run
\end{verbatim}

\subsection{Redox Compatibility}

Of 525 \texttt{.rs} files in the Redox kernel and relibc, Crust lowers items in
about 110 and around 94 of those compile to objects, of which 89 link
together. The upstream code that survives is mostly the memory manager and the
architecture definitions --- \texttt{rmm}'s page tables and flags, the buddy
and bump frame allocators, syscall numbers, layout constants. The
\texttt{heap offset} CrustOS prints at startup is upstream's
\texttt{kernel\_heap\_offset()}, called from code compiled here.

Five of the ninety-four objects cannot join the link, for two reasons that
both follow from compiling files never meant to be one program: Redox ships
one implementation per architecture and Crust flattens module paths, so the
\texttt{aarch64} and \texttt{riscv64} versions of a function collide; and some
files reference symbols from parts of Redox that do not compile.
\texttt{tools/crustos.py} selects a linkable subset with \texttt{nm} and
reports what it dropped and why.

\subsection{Redox Paging}

Redox \texttt{rmm} writes its page flags generically over the architecture, reading
constants off a trait:

\begin{verbatim}
impl<A: Arch> PageFlags<A> {
    pub fn new() -> Self {
        Self::from_data(A::ENTRY_FLAG_DEFAULT_PAGE | A::ENTRY_FLAG_NO_EXEC)
    }
}
\end{verbatim}

CrustOS uses the same shape, and Crust monomorphises it: one implementation of
the flag logic, one instantiation per architecture, every constant resolved at
compile time with no dispatch.

